\documentclass[11pt,a4paper]{article}
\usepackage{amsmath, amssymb, fullpage, mathrsfs, bm, pgf, tikz, bbm}
\usepackage{mathrsfs,amsthm}
\usetikzlibrary{arrows}
\setlength{\textheight}{10in}
%\setlength{\topmargin}{0in}
\setlength{\topmargin}{-0.5in}
\setlength{\parskip}{0.1in}
\setlength{\parindent}{0in}

\begin{document}
	\newcommand{\la}{\leftarrow}
	\newcommand{\lra}{\leftrightarrow}
	\newcommand{\bbN}{{\mathbb N}}
	\newcommand{\bbZ}{{\mathbb Z}}
	\newcommand{\bbQ}{{\mathbb Q}}
	\newcommand{\bbR}{{\mathbb R}}
	\newcommand{\bbC}{{\mathbb C}}
	\newcommand{\bbH}{{\mathbb H}}
	\newcommand{\bbE}{{\mathbb E}}
	\newcommand{\bbP}{{\mathbb P}}
	\newcommand{\dfeq}{\stackrel{\mathrm{def}}{=}}
	\newcommand{\ra}{\rightarrow}
	\newcommand{\Span}{\mathrm{span}}
	\newcommand{\scrP}{\mathscr{P}}
	\newcommand{\rank}{\mathrm{rank}}
	\newcommand{\nullity}{\mathrm{nullity}}
	\newcommand{\Col}{\mathrm{Col}}
	\newcommand{\Row}{\mathrm{Row}}
	\newcommand{\tr}{\mathrm{tr}}
	\newcommand{\ol}{\overline}
	\newcommand{\norm}[1]{||#1||}
	\newcommand{\doubleline}[1]{\underline{\underline{#1}}}
	\newcommand{\elemop}[1]{\stackrel{#1}{\longrightarrow}}
	\newcommand{\Ind}{\mathrm{Ind}}
	\newcommand{\Res}{\mathrm{Res}}
	\newcommand{\End}{\mathrm{End}}
	\newcommand{\cl}{\mathrm{cl}}
	\newcommand{\code}[1]{\texttt{#1}}
	\newcommand\tab[1][0.5cm]{\hspace*{#1}}
	\newcommand{\<}{\langle}
	\renewcommand{\>}{\rangle}
	\newcommand{\qubits}[1]{|{#1}\rangle}
	\newcommand{\powset}{\mathcal{P}}
	\newcommand{\dsum}{\displaystyle\sum}
	\newcommand{\dprod}{\displaystyle\prod}
	
	\newtheorem{lemma}{Lemma}
	
	\section*{Putnam 2025 Solutions}
	
	\section*{Session A}
	\begin{enumerate}
		\item [A1.]
		Let $m_0$ and $n_0$ be distinct positive integers. For every positive integer $k$, define $m_k$ and $n_k$ to be the relatively prime positive integers such that
		\[\frac{m_k}{n_k}=\frac{2m_{k-1}+1}{2n_{k-1}+1}.\]Prove that $2m_k+1$ and $2n_k+1$ are relatively prime for all but finitely many positive integers $k$.
		
		\textbf{Solution.} 
		Let $d_k = m_k - n_k$, which is nonzero. Denote also $c_k = \gcd(2m_{k-1} + 1, 2n_{k-1}+1)$ (which can only be odd). Then we have 
		\[
		d_k = \frac{(2m_{k-1} + 1) - (2n_{k-1} + 1)}{c_k} = \frac{2d_{k - 1}}{c_k}
		\]
		In other words if $o_k$ is the largest odd divisor of $d_k$, then $o_kc_k=o_{k-1}$, or $o_k\mid o_{k-1}$. 
		This means $o_k$ cannot increase and has to be eventually constant (to remain as positive integer). 
		Whenever $o_k=o_{k-1}$ we have $c_k$, which happens for all sufficiently large $k$, as desired. 
		
		\item [A2.] 
		Find the largest real number $a$ and the smallest real number $b$ such that
		\[ax(\pi-x) \le \sin x \le bx(\pi-x)\]for all $x$ in the interval $[0,\pi]$.
		
		\textbf{Answer.} $a=\frac{1}{\pi}; b=\frac{4}{\pi^2}$. 
		
		\textbf{Solution.} We consider each of the bounds separately as follows. 
		
		\textit{Lower bound}. By taking $x\to 0^+$ we know $a=\frac{1}{\pi}$ cannot be improved. On the other hand, $\sin x = \int_0^x \cos t dt$ and $\frac{1}{\pi}x(\pi-x)=\frac{1}{\pi}\int \pi - 2tdt$.
		Considering function $\cos t$ on $[0, \frac{\pi}{2}]$, it is concave, and the line connecting $(0, 1)$ and $(\frac{\pi}{2}, 1)$ is precisely $1 - \frac{2t}{\pi}$.
		Thus $\cos t \ge 1 - \frac{2t}{\pi}$, and so $\sin x \ge \frac{1}{\pi}x(\pi-x)$.
		
		\textit{Upper bound}. By taking $x=\frac{\pi}{2}$ we know $b=\frac{4}{\pi^2}$ cannot be improved. Now let $y = \frac{\pi}{2} - x$ we need to show $\cos y\le 1 - \frac{4y^2}{\pi^2}$, or $1-\cos y\ge \frac{4y^2}{\pi^2}$.
		
		Note the double angle formula gives $1-\cos y=2\sin^2 \frac{y}{2}$, hence it remains to show
		$\sin y \ge 2\sqrt{2}\frac{y}{\pi}$ for $0\le y\le \frac{\pi}{4}$. Note the function $\frac{\sin x}{x}$ decreases in $[0, \frac{\pi}{2}]$ as the derivative $\cos x$ is decreasing, and that $\sin y = 2\sqrt{2}\frac{y}{\pi}$ when $y=\frac{\pi}{4}$, as desired.
		
		\item [A3.] 
		Alice and Bob play a game with a string of $n$ digits, each of which is restricted to be $0$, $1$, or $2$. Initially all the digits are $0$. A legal move is to add or subtract $1$ from one digit to create a new string that has not appeared before. A player with no legal move loses, and the other player wins. Alice goes first, and the players alternate moves. For each $n \geq 1$, determine which player has a strategy that guarantees winning.
		
		\textbf{Answer.} Bob has winning strategy for all $n$. 
		
		\textbf{Solution.} Choose any partition of the $3^n-1$ numbers $\{1, 2, \cdots, 3^n-1\}$ into $\frac{3^n-1}{2}$ pairs such that two numbers in each pair has one digit differing by exactly one, and all other $n-1$ digits being equal. 
		Bob maintains the invariant that after his turn, for each pair, either none of both numbers have been chosen; 
		this is done by choosing the partner (w.r.t. to the pair of numbers) of what Alice has chosen in her round right before Bob. 
		
		Well of course we still need to provide one such parition, one example is to look at the rightmost nonzero digit of each number and flip from 1 to 2 (and vice versa) to form such pair. Below is an example for $n=3$: 
		\[
		(001, 002), (010, 020), (011, 012), (021, 022), 
		(100, 200), (101, 102), (110, 120), (111, 112), (121, 122), 
		\]\[
		(201, 202), (210, 220), (211, 212), (221, 222)
		\]
		Another way to see this is via the following recursive construction: if we have paired $\{1, 2, \cdots, 3^n-1\}$, 
		we may go from $n$-digit to $n+1$-digit via the following: 
		for each pair $(k, \ell)\in \{1, 2, \cdots, 3^n-1\}^2$, we pair $(k, \ell), (3^n+k, 3^n+\ell), (2\cdot 3^n+k, 2\cdot 3^n+\ell)$, and finally pair $(3^n, 2\cdot 3^n)$. 
		
		\item [A5.]
		Let $n$ be an integer with $n \ge 2$. For a sequence $s=(s_1,\dots,s_{n-1})$ where each $s_i=\pm 1$, let $f(s)$ be the number of permutations $(a_1,\dots,a_n)$ of $\{1,2,\dots,n\}$ such that $s_i(a_{i+1}-a_i)>0$ for all $i$. For each $n$, determine the sequences $s$ for which $f(s)$ is maximal.
		
		\textbf{Answer.} There are two such sequences: $s: s_i = (-1)^i$ and $s: s_i =(-1)^{i+1}$. That is, 
		the two sequences such that for all $i < n$, $s_i=-s_{i+1}$. 
		
		\textbf{Solution.} By induction hypothesis, let us assume that this is true for $n-1$ and $n-2$ (the cases where $n=2, 3$ can be checked manually easily). Let $f(s, k)$ be the number of permutations satisfying the condition and where $k=a_n$. Denote $s_{\backslash i}$ as the sequence obtained by removing $s_i$ (and define $s_{\backslash i, j}$ similarly). Then $f(s)=\sum_{k=1}^n f(s, k)$ and each $f(s, k)$ has the following properties:
		
		$$
		f(s, 1)=\boldsymbol{1}_{s_1=-1}\cdot f(s_{\backslash 1})
		\qquad 
		f(s, n)=\boldsymbol{1}_{s_n=1}\cdot f(s_{\backslash n-1})
		$$
		$$
		2\le k\le n-1: f(s, k) =  \binom{n-1}{k-1}\boldsymbol{1}_{s_{k-1}=-1}\cdot \boldsymbol{1}_{s_{k}=1}\cdot f(s_{\backslash k-1, k})
		$$
		We also have $f(s)=f(-s)$ (negation taken elementwise) by mapping $a_1, \cdots, a_n$ to $(n+1-a_1, \cdots, n+1-a_n)$. Therefore we may now compute
		$$
		2f(s) = \sum_{k=1}^n f(s, k) +f(-s, k)=A+B+C
		$$where
		$$
		A= \boldsymbol{1}_{s_1=-1}\cdot f(s_{\backslash 1}) + \boldsymbol{1}_{s_1=1}\cdot f(-s_{\backslash 1})
		$$$$
		B= \boldsymbol{1}_{s_n=1}\cdot f(s_{\backslash n-1}) + \boldsymbol{1}_{s_n=-1}\cdot f(-s_{\backslash n-1})
		$$$$
		C=\sum_{k=2}^{n-1}\left(\boldsymbol{1}_{s_{k-1}=-1}\cdot \boldsymbol{1}_{s_{k}=1}\cdot f(s_{\backslash k-1, k})
		+ \boldsymbol{1}_{s_{k-1}=1}\cdot \boldsymbol{1}_{s_{k}=-1}\cdot f(-s_{\backslash k-1, k})\right)
		$$
		Now, we show that each of $A, B, C$ are maximized when $s$ is altenating. For $A$, exactly one conditions $s_1=-1$ and $s_1=1$ holds for any $s$;
		similarly for $B$; for $C$, at most one condition $s_{k-1}=-1, s_k=1$ and $s_{k-1}=1, s_{k}=-1$ holds; equality when $s_{k-1}=-s_k$.
		Furthermore, if $s$ is alternating, then so are $s_{\backslash 1}, s_{\backslash n-1}$ and $s_{\backslash k, k-1}$.
		This establishes the claim.
		
		To show that the alternating sequences are the only maximizers, note that we need $s_{k-1}=-s_k$ for all $k$ for $C$ to be maximized, which is achieved only when $s$ is alternating.

    \end{enumerate}
    
    \section*{Session B}
    \begin{enumerate}
    	\item [B1.]
    	Suppose that each point in the plane is colored either red or green, subject to the following condition: For every three noncollinear points $A$, $B$, $C$ of the same color, the center of the circle passing through $A$, $B$ and $C$ is also this color. Prove that all points of the plane are the same color.
    	
    	\textbf{Solution.} 
    	Denote red as Type 0 and green as Type 1. The key is to notice that for any point $O$ and circle (circumference) centered at $O$, at most two points on the circumference can be of type different than $O$. 
    	
    	Suppose both types are present; let point $A$ be of type 0 and point $B$ be of type 1. W.l.o.g. let the distance $AB$ to be 1. Now draw two unit circles $\omega_0, \omega_1$ with centers $A$ and $B$ respectively; denote $C, D$ their intersection. 
    	$C, D$ cannot be both type 0 or both type 1, hence w.l.o.g. let $C$ to be of type 0, and $D$ to be of type 1. 
    	In other words only $B, D$ on $\omega_0$ can be of type 1, and $A, C$ on $\omega_1$ can be of type 0. 
    	To bring this logic further: 
    	\begin{itemize}
    		\item For any circle centered at $A$ and radius $r$ with $0<r<2, r\neq 1$, it intersects $\omega_1$ at two points (that are not $A, B, C$ or $D$), hence must be all type 0 except for the two points on $\omega_1$; 
    		\item Opposite conclusion for any circle centered at $B$ and radius $r$ with $0<r<2, r\neq 1$. 
    	\end{itemize}
        However, the region covered by the two bullet points are not mutually exclusive (e.g. the area that is strictly interior of the intersections of $\omega_0$ and $\omega_1$), hence contradiction. 
        
        \item[B2.] 
        Let $f \colon [0,1] \to [0,\infty)$ be strictly increasing and continuous. Let $R$ be the region bounded by $x=0$, $x=1$, $y=0$, and $y=f(x)$. Let $x_1$ be the $x$-coordinate of the centroid of $R$. Let $x_2$ be the $x$-coordinate of the centroid of the solid generated by rotating $R$ around the $x$-axis. Prove that $x_1<x_2$.
        
        \textbf{Solution.} 
        We compute $x_1 = \frac{\int_0^1 xf(x)dx}{\int_0^1 f(x)dx}$, and $x_2=\frac{\int_0^1 xf(x)^2dx}{\int_0^1 f(x)^2dx}$. Thus we shall equivalently show the following: 
        \[
        \int_0^1 xf(x)^2dx\int_0^1 f(x)dx > \int_0^1 xf(x)dx\int_0^1 f(x)^2dx
        \]
        or equivalently (by expanding the integrals)
        \[
        \int_{(x, y)\in [0, 1]^2} xf(x)^2f(y)dxdy > \int_{(x, y)\in [0, 1]^2} xf(x)f(y)^2dxdy 
        \]
        Note that $\int_{(x, y)\in [0, 1]^2} xf(x)^2f(y)dxdy = \int_{0\le x < y\le 1} (xf(x)^2f(y) + yf(y)^2f(x))dxdy$ and 
        $\int_{(x, y)\in [0, 1]^2} xf(x)f(y)^2dxdy = \int_{0\le x < y\le 1}  (xf(x)f(y)^2 + yf(y)f(x)^2)$, 
        allowing us to dissect each such term individually: 
        \[
        xf(x)^2f(y) + yf(y)^2f(x) - (xf(x)f(y)^2 + yf(y)f(x)^2)
        \]\[
        =xf(x)f(y)(f(x)-f(y)) + yf(x)f(y)(f(y)-f(x))
        =(x-y)(f(x)-f(y))f(x)f(y) > 0
        \]
        since $x-y$ and $f(x)-f(y)$ has the same sign (and nonzero) when $x\neq y$. 
        This shows LHS $>$ RHS in the integral, as desired. 
        
        \item [B3.]
        Suppose $S$ is a nonempty set of positive integers with the property that if $n$ is in $S$, then every positive divisor of $2025^n-15^n$ is in $S$. Must $S$ contain all positive integers?
        
        \textbf{Answer.} Yes. 
        
        \textbf{Solution.} Suppose otherwise, let $n$ be the smallest integer not included in $S$. 
        Note that given that $S$ is nonempty, there exists an $m$ such that all positive divisors of $2025^m-15^m\in S$, so $1\in S$. It follows that $n\ge 2$. 
        
        By our assumption, $n\nmid 2025^m - 15^m$ for all $m < n$ and $m\ge 1$. 
        Now write $n=3^a5^b c$ where $c$ is relatively prime to 3 and 5; $c$ is also relatively prime to 2025 and 15. 
        If $c > 1$, take $m=3^a5^b\phi(c)$, we have $n\mid 2025^m - 15^m$. 
        Otherwise, choosing $m=n-1$, we have both $2025^m$ and $15^m$ divisible by $3^{n-1}\cdot 5^{n-1}$, and $n-1\ge a, n-1\ge b$, so $n\mid 2025^m - 15^m$. 
        Both these give contradiction. 
        
        \item [B4.]
        For $n \geq 2$, let $A = [a_{i,j}]_{i,j=1}^n$ be an $n$-by-$n$ matrix of nonnegative integers such that
        \begin{enumerate}
        	\item [(a)] $a_{i,j} = 0$ when $i+j\leq n$;
        	\item [(b)] $a_{i+1,j} \in \{a_{i,j}, a_{i,j}+1\}$ when $1 \leq i \leq n-1$ and $1 \leq j \leq n$; and
        	\item [(c)] $a_{i,j+1} \in \{a_{i,j}, a_{i,j}+1\}$ when $1 \leq i \leq n$ and $1 \leq j \leq n-1$.
        \end{enumerate}
        Let $S$ be the sum of the entries of $A$, and let $N$ be the number of nonzero entries of $A$. Prove that
        \[S \leq \frac{(n+2)N}{3}.\]
        
        \textbf{Solution.} 
        Call the $n\times n$ matrix satisfying the problem condition $n$-\textit{nice}; 
        note that $B$ defined as $B=A[2:n, 1:n-1]$ is $n-1$-nice. 
        Correspondingly, we denote $S'$ and $N'$ as the sum of entries and number of nonzero entries of $B$, respectively. 
        By inductive hypothesis we have $S'\le \frac{(n+1)N'}{3}$. 
        
        Denote $a = \sum_{i=1}^n a_{i, n}$; we have $S=S'+a$. 
        Denote $b = \sum_{i=1}^n \boldsymbol{1}_{a_{i, n} > 0}$; we have $N=N'+b$. 
        With $a_{i, n}-a_{i-1, n}\in \{0, 1\}$ for all $i\ge 2$, the $k$-th smallest nonzero entry of $a_{i, n}$ (for $k\le b$) must be at most $k$, 
        hence $a\le \frac{b(b+1)}{2}$. 
        Meanwhile, $a_{i, n} - a_{i, n - j}\le j$ so if $a_{i, n} > 0$, 
        the number of nonzero entries among $a_{i, 1}, \cdots, a_{i, n}$ is at least $a_{i, n}$. 
        It then follows that $A\ge \sum_{i=1}^n a_{i, n} = a$. 
        Hence $a\le \min \{N, \frac{b(b+1)}{2}\}\le \min \{N, \frac{b(n+1)}{2}\}$ 
        (since $b\le n$). 
        It then follows that 
        \[
        S = S' + a \le \frac{(n+1)N'}{3} + a
        \le \frac{(n+1)N'}{3}  + \frac{N}{3} + \frac{b(b+1)}{3}
        \le \frac{(n+1)N'}{3}  + \frac{N}{3} + \frac{b(n+1)}{3}
        =\frac{(n+2)N}{3} 
        \]
        using $N=N'+b$, as desired. 
        
        \item [B5.]
        Let $p$ be a prime number greater than $3$. For each $k \in \{1,\dots,p-1\}$, let $I(k) \in \{1,2,\dots,p-1\}$ be such that $k \cdot I(k) \equiv 1 \pmod{p}$. Prove that the number of integers $k \in \{1,\dots,p-2\}$ such that $I(k+1)<I(k)$ is greater than $p/4-1$.
        
        \textbf{Solution.}
        Consider the number $r(k)\in \{0, \cdots, p-1\}$ such that $r(k)\equiv I(k+1)-I(k)$. Note that $I(p-1)=p-1$ and $I(1)=1$, and $r(k)$ is either $I(k+1)-I(k)$ or $I(k+1)-I(k)-p$. Thus the number of interest is $\frac{1}{p}\left(\sum_{k=1}^{p-2} r(k)-(p-2)\right)$.
        
        Now, $r(k)\equiv \frac{1}{k+1}-\frac{1}{k}=\frac{-1}{k(k+1)}$, so $k(k+1)r(k)\equiv -1$. We now note the following:
        
        - $k(k+1)\equiv \ell(\ell+1)$ if and only if $k=\ell$ or $k+\ell\equiv p-1$; in other words $r(1), \cdots, r(\frac{p-1}{2})$ are all different.
        
        - $I(\frac{p+1}{2})=2$ and $I(\frac{p-1}{2})=p-2$, so $r(\frac{p-1}{2})=4$.
        
        Thus $r(1), \cdots, r(\frac{p-3}{2})$ are all distinct, nonzero, and not equal to 4. Note that $r(p-2), \cdots, r(\frac{p+1}{2})$ are equal to these too. Therefore we may lower bound the sum as
        
        $\left(\sum_{k=1}^{p-2} r(k)\right)\ge 2(1 + 2 + \cdots + \frac{p-1}{2}) - 4 = \frac{p^2-1}{4}-4$.
        
        For $p\ge 13$, we have $\frac{p^2-1}{4}-4\ > \frac{p(p-1)}{4}+2$, and so $\frac {1}{p}\left(\sum_{k=1}^{p-2} r(k) -(p-2)\right) > \frac{p-1}{4} - 1$. The LHS of this inequality must be an integer, giving the bound $\frac{p-1}{4}$ for $p\equiv 1\pmod{4}$ and $\frac{p-3}{4}$ for $p\equiv 3\pmod{4}$, as desired.
        
        It remains to check $p=5, 7, 11$. For all cases we have $I(\frac{p+1}{2})=2 < p-2=I(\frac{p-1}{2})$. In addition for $p=11$ we have $I(3)=4<6=I(2)$ and $I(9)=5<7=I(8)$.
    \end{enumerate}
\end{document}